\ssr{ЗАКЛЮЧЕНИЕ}

Не существует универсальной системы управления реестром Docker-образов, которая подходила бы для всех сценариев использования. Разные решения имеют свои сильные и слабые стороны, поэтому выбор конкретной системы определяется поставленными задачами и требованиями организации.

Как показало исследование, для корпоративных сред важны такие функции как безопасность, отказоустойчивость и интеграция с существующей инфраструктурой. В таких случаях можно выбирать из решений уровня Harbor, Quay или облачных сервисов крупных провайдеров. Однако для тестовых сред или небольших проектов могут подойти и более простые решения.

В данной научно-исследовательской работе был:
\begin{enumerate}
	\item проведен анализ предметной области управления реестрами Docker-образов;
	\item выполнена классификация задач и требований к системам управления реестрами;
	\item проведен сравнительный анализ 30 систем по сформулированным критериям;
	\item определены преимущества и ограничения рассмотренных решений.
\end{enumerate}

В рамках поставленной работы были выполнены все задачи. Цель работы — проведение сравнительного анализа систем управления реестрами на основе платформы Harbor — достигнута.

На основе проведенного исследования можно сделать вывод, что система Harbor является одним из наиболее сбалансированных решений среди открытых систем и может быть рекомендована для организаций, ищущих полнофункциональную платформу без привязки к конкретному облачному провайдеру. Однако окончательный выбор системы должен учитывать конкретные требования и контекст использования.

